\documentclass[11pt]{article}
\usepackage[margin=1in, headheight=14pt]{geometry}
\usepackage{amsmath, amssymb}
\usepackage{hyperref}
\usepackage{enumitem}
\usepackage{algorithm}
\usepackage{algpseudocode}
\usepackage{fancyhdr}
\usepackage{tcolorbox}
\usepackage{microtype}
\usepackage{tikz}
\usepackage{pgfplots}
\pgfplotsset{compat=1.18}
\usetikzlibrary{arrows.meta, positioning, calc, shapes.geometric}

\input{notation.tex}
\input{zk-notation.tex}
\input{environments.tex}
\input{lecture-macros.tex}

\pagestyle{fancy}
\fancyhf{}
\lhead{EECS 498/CSE 598: Zero-Knowledge Proofs}
\rhead{Winter 2026}
\lfoot{Lecture 14: Frontends, Transcripts, and ASIC-Model Frontends}
\rfoot{\thepage}
\renewcommand{\headrulewidth}{0.4pt}
\renewcommand{\footrulewidth}{0.4pt}

% Lecture-specific macros
\tcbuselibrary{skins, breakable}

\newtcolorbox{highlight}[1][]{
  colback=blue!5!white,
  colframe=blue!60!black,
  fonttitle=\bfseries,
  #1
}

\begin{document}

%-----------------------------------------------------------------------
% Title Block
%-----------------------------------------------------------------------
\begin{center}
  {\LARGE \textbf{EECS 498 / CSE 598: Zero-Knowledge Proofs}}\\[6pt]
  {\large Winter 2026}\\[4pt]
  {\large \textbf{Lecture 14: Frontends, Transcripts, and ASIC-Model Frontends}}\\[4pt]
  {\large Instructor: Paul Grubbs \quad
    \href{mailto:paulgrub@umich.edu}{\texttt{paulgrub@umich.edu}} \quad
    Beyster 4709}
\end{center}

\vspace{0.5em}
\hrule
\vspace{1em}

%-----------------------------------------------------------------------
% Agenda
%-----------------------------------------------------------------------
\section*{Agenda}
\begin{itemize}
  \item Frontends and witness extension
  \item Computation transcripts
  \item Translation soundness and nondeterminism
  \item ASIC and CPU programming models
  \item ASIC-model frontend constraints
\end{itemize}

\hrule
\vspace{1em}

%-----------------------------------------------------------------------
\section{Frontends}
%-----------------------------------------------------------------------

A frontend is the layer that connects an ordinary programming abstraction to
the algebraic representation consumed by a proof system.

\begin{highlight}[title={Frontend Output Triple}]
For a source program
\[
  P(x,w) \in \{0,1\},
\]
a frontend produces:
\begin{enumerate}
  \item an algebraic representation $C$ of the computation,
  \item a witness-extension procedure $E$ that maps $(x,w)$ to compiled inputs
        $(x',w')$,
  \item a correctness guarantee of the form
        \[
          P(x,w)=1
          \iff
          C(x',w')=1
          \qquad\text{for } (x',w') = E(x,w).
        \]
\end{enumerate}
\end{highlight}

The point of the source program $P$ is not merely syntactic convenience.  It
captures the intended relation using familiar control-flow and data-abstraction
mechanisms.  The compiled representation $C$ is then the object actually
consumed by the backend proof system.

\begin{definition}[Witness extension]
The witness-extension procedure $E$ is the auxiliary computation that turns a
source-level witness into the full compiled witness required by the algebraic
representation.
\end{definition}

The frontend therefore does more than compilation in the ordinary programming
language sense.  It also specifies how the hidden witness must be enriched so
that the backend relation can be checked.

%-----------------------------------------------------------------------
\section{Programs and Transcripts}
%-----------------------------------------------------------------------

The central witness object for a frontend is often a \emph{transcript} of the
execution.  A transcript records the values of intermediate variables rather
than only the final output.

\begin{definition}[Computation transcript]
For a deterministic program execution, a transcript is the ordered collection
of internal states or intermediate variable values produced during that
execution.
\end{definition}

Consider the source program
\[
\begin{array}{l}
\algo{factorial}(n): \\
\quad \algo{acc} \leftarrow 1 \\
\quad \text{for } i = n,n-1,\ldots,1: \\
\qquad \algo{acc} \leftarrow \algo{acc} \cdot i \\
\quad \text{return } \algo{acc}.
\end{array}
\]
For $n=4$, one possible transcript is
\[
  (\algo{acc}_0,i_0) = (1,4),\ 
  (\algo{acc}_1,i_1) = (4,3),\ 
  (\algo{acc}_2,i_2) = (12,2),\ 
  (\algo{acc}_3,i_3) = (24,1).
\]

\begin{remark}
The proof objective is transcript verification rather than transcript
generation.  The prover supplies the internal values, and the verifier checks
that adjacent states satisfy the program relation.
\end{remark}

This point is especially important for zero knowledge.  The prover can provide
auxiliary values that would be inefficient or unnatural to recompute inside the
constraint system, as long as the constraints verify their consistency.

%-----------------------------------------------------------------------
\section{Structural Design Issues}
%-----------------------------------------------------------------------

\subsection{Nondeterminism}

Frontend design usually treats proving as a nondeterministic verification task.
An expensive computation can be replaced by a guessed value together with a
cheap consistency check.  Multiplicative inverses are the standard example:
instead of computing $a^{-1}$ inside the circuit, the prover supplies a value
$b$ and the circuit checks
\[
  ab = 1.
\]

\subsection{Enriching the Target Model}

The target algebraic model need not be plain $\algo{R1CS}$.  Modern backends
often expose richer primitives such as lookup arguments, custom gates, or
variants of $\algo{R1CS}$ with extra structure.  A frontend can exploit these
features if doing so yields a more compact or more natural representation of
the source computation.

\subsection{Translation Soundness and Completeness}

The frontend translation has two basic correctness obligations:
\begin{itemize}
  \item \textbf{Completeness:} every valid source-level execution should map to
        a valid compiled witness.
  \item \textbf{Soundness:} every accepting compiled witness should correspond
        to a valid source-level execution of the intended program.
\end{itemize}

The first obligation is ordinary compiler correctness in the proving
direction.  The second obligation is stronger: it rules out bugs in the
translation that accidentally enlarge the accepted relation.

%-----------------------------------------------------------------------
\section{ASIC and CPU Programming Models}
%-----------------------------------------------------------------------

\begin{definition}[ASIC-model frontend]
An ASIC-model frontend translates the source program into a fixed algebraic
computation by expanding the executed operations line by line into the target
constraint system.
\end{definition}

For a loop such as
\[
\begin{array}{l}
i \leftarrow 0; \\
\text{for } j = 0; j < 10; j++ \text{ do } i \leftarrow i+1,
\end{array}
\]
the ASIC-model translation unrolls the loop into static single-assignment
(\algo{SSA}) form:
\[
  z = 0,\quad
  z_0 = z + 1,\quad
  z_1 = z_0 + 1,\quad
  \ldots,\quad
  z_9 = z_8 + 1.
\]

The fixed circuit structure is the defining feature.  Every program path that
may occur must already be reflected in the compiled algebraic object.

By contrast, a CPU-model frontend compiles the program to an instruction set
architecture and reasons about the resulting execution trace: opcode sequence,
register states, and memory accesses.  Mutable memory and data-dependent
control flow are then handled through transcript constraints rather than by
fully unrolling the source program into one static arithmetic object.

\begin{remark}
The CPU model is closer to ordinary compilation pipelines because the frontend
can target assembly or bytecode rather than a bespoke arithmetic intermediate
representation.
\end{remark}

%-----------------------------------------------------------------------
\section{ASIC-Model Frontends in Practice}
%-----------------------------------------------------------------------

Representative ASIC-model frontends include
\[
  \text{\algo{Buffet}, \algo{Pantry}, \algo{Pequin}, \algo{Pinocchio},
  \algo{xjSnark}, \algo{CirC}, \algo{Noir}, \algo{Circom},
  \algo{Lurk}, \algo{Bellman}}.
\]

One important structural distinction is between embedded domain-specific
languages and compiled languages:
\begin{itemize}
  \item embedded languages reuse a host language and embed proving constructs
        into it;
  \item compiled languages define their own syntax and type system before
        lowering to the backend model.
\end{itemize}

\algo{P4} is an example of the embedded-language side of this distinction.

Several restrictions appear repeatedly in ASIC-model frontends:
\begin{enumerate}
  \item the dominant scalar datatype is usually the backend field element;
  \item loops require static bounds so that the computation can be unrolled;
  \item some source-level computations are enforced inside the proof, while
        others are left unchecked unless the programmer inserts explicit
        constraints.
\end{enumerate}

These restrictions are not accidental.  They reflect the fact that the target
object is still a fixed algebraic relation.  The frontend succeeds when the
source language exposes enough ordinary programming structure to be usable while
still compiling cleanly into that fixed target.

\end{document}
